% =====================================================================
% Apêndice A — resumo do plano de análise e dos desvios
% Fonte: docs/pre-especificacao.md (fechada em 2026-09-21, commit c386514).
% =====================================================================

\chapter{Plano de análise do Ensaio 1 e desvios registrados}
\label{ap:preespecificacao}

O plano de análise foi fechado e versionado em 21 de setembro de 2026. Ele não
é um pré-registro em sentido estrito: foi escrito depois dos diagnósticos de
desenho (distribuição da dose, dispersão das tendências municipais, efeito
mínimo detectável, relação entre aptidão e dose) e antes de qualquer estimativa
de efeito. A fronteira que ele protege é entre desenho e efeito. O
Quadro~\ref{qua:plano} resume as escolhas, e o Quadro~\ref{qua:desvios}, os
desvios posteriores.

\begin{quadro}[htb]
\centering
\caption{Escolhas do plano de análise}
\label{qua:plano}
\small
\begin{tabular}{|p{4.2cm}|p{10cm}|}
\hline
\textbf{Item} & \textbf{Escolha} \\
\hline
Desfecho primário & Peso médio ao nascer, por município e mês \\
\hline
Exposição primária & Fração da gestação posterior à vigência da lei \\
\hline
Cultura-âncora & Banana, por ter área positiva em 169 municípios e índice de
Gini da dose de 0,86 \\
\hline
Janela da dose & 2015--2018; 2010--2014 como sensibilidade \\
\hline
Grupo de comparação & Municípios de dose nula, definidos por baixa aptidão
agroclimática; área nula como sensibilidade \\
\hline
Cortes & Período anterior até dez/2018 e posterior a partir de out/2019;
coortes de transição descartadas (9,4\% das células) \\
\hline
Critério de falsificação & Piso de efeito esperado de 15 g \\
\hline
Família confirmatória & Peso médio e óbito fetal, com correção de Holm \\
\hline
Exploratório declarado & Exposição por trimestre; baixo peso e prematuridade;
canal de intoxicação; janela 2010--2014; definições alternativas do grupo de
comparação \\
\hline
Falsificação & Notificações de intoxicação intencional, que não devem se mover \\
\hline
Vedações & Não trocar o desfecho primário; não descer de unidade geográfica;
não reportar como achado coeficiente abaixo do efeito mínimo detectável; não
usar o canal hídrico nem o teste de direção do vento \\
\hline
\end{tabular}
\legend{Fonte: elaboração própria, a partir do plano de análise versionado.}
\end{quadro}

\begin{quadro}[htb]
\centering
\caption{Desvios do plano de análise}
\label{qua:desvios}
\small
\begin{tabular}{|p{2.1cm}|p{5.2cm}|p{6.9cm}|}
\hline
\textbf{Data} & \textbf{O que mudou} & \textbf{Motivo} \\
\hline
21/09/2026 & Grupo de comparação primário passa a ser o de baixa aptidão
  & O índice FAO-GAEZ foi obtido para os 184 municípios \\
\hline
22/09/2026 & Parâmetro principal passa da curva, $ACR(d)$, ao nível, $\ATT$
  & A derivada agregada é dominada pela faixa de dose em que o erro-padrão
    supera a estimativa \\
\hline
22/09/2026 & A hipótese de tendências paralelas fortes deixa de ser descrita
  como não testável
  & A verificação anterior ao tratamento não a isola, mas pode falsificá-la;
    foi implementada \\
\hline
22/09/2026 & A correção de Holm passa a ser calculada
  & Estava declarada e não implementada \\
\hline
\end{tabular}
\legend{Fonte: elaboração própria, a partir do plano de análise versionado.}
\end{quadro}

As decisões marcadas ao longo do texto como dependentes do orientador, se
adotadas, entrarão neste quadro com data e motivo, e os resultados originais
serão preservados.
